\documentclass[11pt]{article}
\usepackage[letterpaper,margin=.85in]{geometry}
\usepackage[T1]{fontenc}
\usepackage{lmodern,microtype,amsmath,amssymb,booktabs,array}
\usepackage{xcolor,enumitem,fancyhdr}
\usepackage[hyphens]{url}
\usepackage{hyperref}
\definecolor{navy}{HTML}{17324D}
\definecolor{teal}{HTML}{176B73}
\hypersetup{colorlinks=true,linkcolor=navy,urlcolor=teal,citecolor=teal,
 pdftitle={Hierarchical ORAM: Communication, Cloud Cost, and Latency}}
\setlength{\parindent}{0pt}\setlength{\parskip}{5pt}
\setlength{\emergencystretch}{2em}\setlength{\headheight}{14pt}
\setlist{nosep,leftmargin=1.6em}
\pagestyle{fancy}\fancyhf{}
\fancyhead[L]{\small\textcolor{navy}{Hierarchical ORAM: cloud deployment}}
\fancyhead[R]{\small Current construction and deployment}
\fancyfoot[C]{\thepage}
\begin{document}
{\color{teal}\large TECHNICAL REPORT\par}
{\color{navy}\LARGE\bfseries Hierarchical ORAM:\\Communication, Cloud Cost, and Latency\par}
17 October 2026\quad$\cdot$\quad Single projected AES per lookup

\section*{Result}
%%RESULTPROSE%%

%%SUMMARY%%

These are newly sized and simulated deployments of the same construction:
one address PRF per occupied table, private random dummy addresses,
short-output AES pruning in both lookup and rebuilding, 64-entry lookup
buckets, and Boolean fingerprints with packed payloads. All resource and
price tables below use this protocol. There is no separate dummy PRF or
dummy epoch key. The budget charges both parties and all seven stored
states; each query runs against one ready state.

When a few hundred milliseconds are acceptable, $B=2048$ is the cheaper
tested choice. $B=1024$ provides more latency margin. Rebuilding and
prefetching remain the larger worker pool. The predictions assume that
loading/transferring circuit material dominates CPU computation. They are
conditional analytical results and finite deterministic traces, not AWS
benchmarks or latency guarantees. All MB/GB/TB are decimal; MiB/GiB are binary.

For parties on different cloud platforms, Section~\ref{sec:cross-cloud}
extends the analysis to private AWS--Google links and larger batches.
The resulting budgets include cross-cloud bandwidth and allow a few
seconds of query latency.

\section{The current construction}
\subsection{Tables, reads, and representations}
The branching factor is two. Logical table capacities are
$n=B,2B,\ldots,N$; each table has $2n$ main slots, grouped into
$m_n=2n/64$ lookup buckets, plus one global stash $G(n)$. The base is a
hash table of capacity $B$, not a dense stash. A lower table is occupied
half of a binary-counter cycle; the largest table is always occupied.
With $L=\log_2(N/B)$, the mean number of occupied tables is $A=L/2+1$.

Hash each application key once to a secret 128-bit canonical fingerprint
using a PRF. The numerical workload assumes a fixed 32-byte application
key, charged as two AES-CMAC encryptions and its shared input \cite{nist-cmac}. A bitonic
sort deduplicates the batch once; saved private comparison controls support
final result propagation and inverse routing. There is no recursive
position map and no repetition of request deduplication at hierarchy levels.

Visit occupied tables from small to large. In each, reveal one selected
bucket address and privately scan that bucket and the global stash. A
match clears the entry's occupancy flag. Once a request has matched, later
tables receive dummy probes. Reinsert one representative per deduplicated
request, including a keyed absent tombstone when appropriate. The retained
query fingerprint supplies the reinsertion key, so a second copy need not
be extracted from the matching entry. Shared buckets and stashes receive
one combined batch update, not conflicting per-query writes.

Each entry holds a 128-bit fingerprint and two Boolean flags (occupied and
present), with one 16-byte label per bit per party. Its 256-bit payload is
a 32-byte XOR share per party. This is an explicit word-OT representation;
two ordinary GC labels do not encode an arbitrary 256-bit value. The
persistent per-party entry size is
\[
 E=16(128+2)+32=2112\text{ bytes}.
\]
We use 32 bytes per half-gate AND and 160 bytes per packed conditional swap
(80 in each direction). A Boolean conditional swap costs one AND per
moved bit. Therefore a record swap costs
\[
 \underbrace{4096}_{128\text{-bit key}}+
 \underbrace{64}_{\text{two flags}}+
 \underbrace{160}_{\text{payload}}=4320\text{ bytes},
\]
plus 32 bytes per retained routing bit. The persistent fingerprint remains
Boolean throughout this costed construction, including rebuilding.

Fingerprint equality uses 127 ANDs per scanned slot. Match, occupancy,
active-request, and present handling have a 134-AND/slot allowance in the
ledger, including equality. The allowance is conservative over a 130-AND
core using one-hot XOR reductions. Packed extraction is charged separately
at 160 bytes per candidate slot. Every lookup still scans a bucket and
stash at every occupied table, including after a hit.

\subsection{Rebuilding and recursive refinement}
A carry rebuilds capacity $n$ once per $2n$ requests; the largest table
rebuilds once per $N$. Rebuilds are blocking for that state. Availability
comes from other independent states, without a deamortized ORAM algorithm.
The stages, including data movement, are:
\begin{enumerate}
\item Compact source main arrays. At the largest rebuild, discard absent
tombstones. Append old stash entries and compact again to readmit them.
\item Retain a public $n$ positions and compute fresh secret routing tags
under the destination table's new epoch key. Dummy positions still receive
fixed-shape work. Spread the compacted slots into half-full 512-slot buckets.
\item Run the binary 512-slot butterfly. A MergeSplit planner consumes only
one routing bit per slot, but its secret controls move the complete entry
and all unconsumed routing bits through ORCompact switches.
\item Refine each 512-slot parent into eight 64-slot leaves. Plan per-leaf
membership, identify excess entries, collect them once into a 64-slot local
stash, then recursively partition the feasible remainder. A public retained
mask is inverse-routed to identify which original slots remain in each leaf.
\item Compact all local residual stashes into the table's global stash,
with a fixed public capacity and a charged overflow check.
\end{enumerate}
For a power-of-two $q$-slot ORCompact, the ledger uses $q\log_2q/2$ swaps
and a constructive planner allowance $q\log_2q-q/2-\log_2q+1$ ANDs. These
are the project's cost derivation, not gate counts quoted from the paper.
The routing method follows the supplied Microrama and ORCompact sources
\cite{construction,orcompact}.

Global stash capacities are selected using a Chernoff/Poisson residual
bound and rounded to multiples of eight. At $B=1024$, the base stash is
56 slots and the largest stashes are 112/264 slots for $N=2^{28}/2^{32}$.
These capacities and the overflow checks are included in read and rebuild
costs. Their existing failure accounting uses $Q=2^{32}$ accesses;
extending the failure budget to a full year remains a stated validation
task, rather than an implicit annual guarantee.

\section{One PRF and private dummy addresses}
For every visit, each party samples a fresh private $t_n$-bit share,
where $t_n=\log_2(2n/64)$. Inside the circuit, form
$r=r_G\oplus r_E$ and reveal only
\[
 j=\operatorname{mux}(\mathrm{active},H_{K_{n,e}}(x),r).
\]
The dummy address is not opened separately. If both parties already knew
it, comparing it with $j$ would identify dummy visits. Generate and input
shares for every visit, regardless of the hidden active flag. Our shared
input allowance is $64t_n$ bytes: $48t_n$ forward and $16t_n$ reverse.
The $t_n$-AND address mux is separately included in lookup logic.

The shares and their input labels can be prepared with the next circuit
batch, so no extra dependent online coin-toss round is assumed. Prefetch
accounts for both traffic directions and conservatively reserves the full
input transcript in RAM. It does not treat that transcript as garbled
gate bytes. The real-address PRFs can be computed before previous-table
hit bits are available; only output selection depends on those bits.

The address PRF is full-round AES-128 followed by a fixed public output
projection. Keep output bytes $(0,13,10,7)$, concatenated in that order,
and take the first $t$ bits in a fixed bit order. These are zero-based byte
positions in AES column-major serialization. Inverting the final
ShiftRows maps the selected bytes to one round-9 column. For $t\leq32$,
all 16 S-boxes in rounds 1--8, four in round 9, and $\lceil t/8\rceil$
in round 10 suffice. Using the same 32-AND S-box as the full 5120-AND
datapath gives the upper bound
\[
 g(t)=32\left(128+4+\left\lceil t/8\right\rceil\right).
\]
MixColumns and AddRoundKey are linear. No AES round is removed and both
lookup and rebuilding use the identical projection
\cite{nist-fips197,nist-aes}. We charge a complete S-box even for a partly
used final byte. A completed full-output circuit would still cost 5120
ANDs; this saving requires pruning the unused computations.

%%PRUNING%%

Canonical fingerprints remain 128 bits. The two initial CMAC block
encryptions remain full-output AES, and each table epoch receives a fresh
independent secret address key with a full 1280-AND key schedule and shared
key input. There is no separate dummy key. For weights $w_n=1/2$ at lower
tables and $w_N=1$, the amortized PRF datapath charge is
\[
 C_{\rm PRF}=327680+64\sum_n w_n g(t_n),
\]
because the lookup weight and $n/\mathrm{rebuildInterval}(n)$ both equal
$w_n$. Private dummy inputs add $64\sum_nw_nt_n$ bytes/query. Epoch-key
setup and inputs add $49152/B$ bytes/query: 48 at $B=1024$, 24 at $B=2048$.

Hash reuse already occurs within one table: one projected output supplies
all butterfly and refinement address bits. Reusing address prefixes across
tables or epochs can expose correlations and the hidden real-to-dummy
transition; independent contexts are retained. Carrying an extra 128-bit
Boolean hash through every swap costs another 4096 bytes/swap. One saved
projected AES pays for only 33.25--34 such extra swaps. This makes that
straightforward full-tag retention strategy unattractive, without ruling
out a separately proved packed or selectively retained representation.

\section{Communication breakdown and circuit interfaces}
%%COMMUNICATION%%

Rebuilding contributes 10.883/14.662 MB/query at $B=1024$, more than half
the useful total. Boolean fingerprint movement alone contributes
7.911/10.685 MB/query. Payload movement is small by comparison. The next
table exposes the rebuild stages so that source compaction, recursive
splitting, stash readmission, and residual collection are all visible.

%%STAGES%%

The cost requires compatible circuit composition: a common session-wide
Free-XOR offset with odd color bit, unique gate/wire namespaces, and direct
import of committed state-wire pairs. Immutable fingerprints and expanded
keys may feed fresh circuits. Mutable occupancy outputs receive fresh
wire pairs; changed values must never reuse the old pair. These are
requirements on the eventual compiler and implementation.

Prefetched query ports also need fresh pairs. With common offset $\Delta$,
a translation from $W^b=W^0\oplus b\Delta$ to $V^b$ uses one public
16-byte value $T=W^0\oplus V^0$. Reserve two 128-wire ports per occupied
table, totaling 4096 bytes: one for the PRF/stash lane and one for the
selected bucket equality lane. Cold stages without a bucket cache need
at most the first port; retaining both is conservative. Every cached
version and every query translation is single-use.

If instead each kernel independently re-encodes all persistent inputs,
generic two-row unary bridges add approximately 5.849/8.279 MB/query at
$B=1024$, before further internal boundaries. The priced construction uses
direct state composition; this overhead cannot simply be ignored by an
implementation using unrelated labels. Cache binding and batch finalization
remain compilation/security-composition TODOs.

\section{Caching and the AWS resource bottleneck}
\subsection{Small-table RAM caches and lookahead}
Cache the small tables near the batch-sized base, through capacity $2^{16}$.
Queue fresh immutable-fingerprint equality circuits and keep their state
in RAM. Keep each table's global stash state and one future batch of its
stash, PRF, and private-randomness material in RAM at every level. The
remaining occupancy/present/active finalizer is generated dynamically.

A table of logical capacity $n$ lasts $n$ requests while present and has
$2n/64$ buckets, so a bucket receives 32 accesses on average over its
lifetime. That is not a useful per-batch upper bound. At the base with
$B=1024$, there are only 32 buckets and a mean of 32 accesses per bucket
in that single batch. A four-version queue covers at most 12.5\% of those
requests. A larger adaptive queue is needed near the base.

With $X\sim\operatorname{Binomial}(B,1/m_n)$, choose the smallest queue
width $w$ satisfying $m_n\Pr[X>w]$ below the table's share of a $10^{-3}$
batch any-miss budget. At $B=1024$, the seven widths are
$[59,37,24,17,12,9,7]$. This is a performance budget under the ideal-PRF
independent-address model, not an ORAM failure probability. Duplicates are
removed before real probes; padded inactive requests get independent
private dummy addresses.

%%CACHE%%

The RAM table conservatively reserves the full 134-AND bucket allowance;
the prefetched immutable equality component uses 127 ANDs/slot. Useful
prefetched bytes remain in total communication. Add only unused versions
at table retirement and input bridges. Skip the final refill before a
publicly known retirement. Cold buckets load state on demand and receive
fresh circuits streamed to RAM; no cold-bucket GC disk queues are part of
the priced deployment.

\subsection{Network, EBS, and memory accounting}
Use \texttt{c6in.32xlarge}, with 256 GiB RAM, 200 Gbps network per direction,
and 100 Gbps aggregate EBS bandwidth \cite{aws-compute,aws-network}. The
modeled active rates per worker are
\begin{align*}
 C_{\rm NIC}&=25\cdot0.65\cdot0.70=11.375\ {\rm GB/s},\\
 C_{\rm EBS}&=12.5\cdot0.65=8.125\ {\rm GB/s},\\
 C_{\rm DRAM}&=56\ {\rm GB/s}.
\end{align*}
The NIC factors are network and distributed efficiency. The DRAM rate is
an explicit effective-loading proxy, not an AWS specification. Same-AZ pool sizing
then limits modeled mean utilization to 70\%; this factor is not applied
again to active service durations. All traffic receives 5\% framing.

Charge half the modeled external record read/write boundary volume as
intra-party sends and receives in each party's limiting NIC direction.
This assumes a distributed storage owner sends records to workers; gp3
is not assumed to support cross-host multi-attach. Actual gather/scatter
may differ. If this internal volume is $I$ per party per direction, both
parties together send $2I$, not $4I$: counting each matching receive again
would double-count network traffic.

%%NETWORK%%

If $F$ is cross-party forward traffic, the limiting directional capacity
uses $1.05(F+I)$, not the all-network sum. The reverse direction includes
packed OTs and private-randomness input traffic and is checked separately.
State I/O and DRAM loads are not extra 2PC communication, although remotely
transferring a record is network traffic as represented by $I$.

%%CAPACITY%%

Retain each 1024-slot MergeSplit and 512-slot refinement kernel in a small
RAM tile. Materialize only the outer butterfly boundaries; source
compactions, stash readmission, and global residual collection still incur
their charged large-array passes. A conservative separate source/final
boundary pass is also included. Keeping inner switches local is essential:

%%IO%%

The complete state does not fit in the small-table cache: it is
2.268/36.284 TB per party per state at $B=1024$. With tiled movement and
fresh GC streaming, network transfer remains the main aggregate throughput
limit. Counting two GC-memory passes, full switch-level record movement,
and source/final boundary traffic, no rebuilt level in the tested cases
requires more than 32.92 GB/s per worker to keep up with its modeled
NIC/EBS service. This is below the 56 GB/s DRAM assumption. EBS or DRAM
can still determine individual read-stage service.
Spilling every switch layer or writing/rereading all GCs changes the
bottleneck. Planner labels, OT matrices, state update amplification, and
cache-owner distribution require implementation measurements.

\section{Fleet scheduling and annual cost}
\subsection{Worker pools and independent states}
There are $F$ foreground and $R$ background workers per party. Each pool
jointly serves one job at a time. The foreground pool handles a batch;
the background pool performs that state's rebuild and prepares its next
circuit batch. Background jobs use preemptive shortest remaining work,
allowing a small rebuild of another state to interrupt a large rebuild.
A state stays unavailable until its whole rebuild and prefill finish.
This changes scheduling across states, not the within-state ORAM algorithm.

States start warm at evenly spaced valid public counter offsets. Dispatch
to the oldest ready state, queueing if none is available. The model uses a
static lookup database: independent states contain the same logical data
but evolve their own ORAM placements. Each request changes one state and
is charged once. Application updates that must reach every state require
an additional consistency protocol and are outside this workload.

For the same-AZ deployment, pool sizes are chosen separately so mean foreground and background
occupancy, including modeled round/load waits, stay at or below 70\%.
The full binary-counter trace is then simulated. Each run spans two global
$N/B$ batch cycles: about 14.91 hours at $N=2^{28}$ or 9.94 days at
$N=2^{32}$. The admission queue is included in the reported latency.

Seven states are used for the headline results. Five states are a useful
smaller deployment comparison, using exactly the same current protocol
and worker pools; they still exhibit occasional multi-second stalls in
the larger-database traces. Additional states cost storage and cache
capacity, not seven copies of the worker fleet:

%%STATES%%

The maxima are observations over two largest-table rebuild starts per
trace. They are neither worst-case proofs nor annual tail guarantees.
Cache misses, failures, bursty arrivals, and real checkpoint/suspension
costs are not sampled. The preemption and storage models are assumptions
to validate in a prototype.

\subsection{AWS prices and storage reservations}
Use the official us-east-1 price snapshot published 15 October 2026 and
checked 17 October: Linux On-Demand \texttt{c6in.32xlarge} at \$7.2576/hour;
gp3 capacity \$0.08/GiB-month; added throughput \$0.04/(MiB/s)-month; and
added IOPS \$0.005/IOPS-month \cite{aws-prices,aws-gp3}. The feed's
throughput unit is GiB/s-month and is divided by 1024. Annualization uses
8760 instance hours and 12 storage-billing months.

Each worker receives eight gp3 volumes, each provisioned at 1491 MiB/s
and 23,842 IOPS. This buys 12.507 GB/s aggregate raw throughput per worker,
enough for the 12.5 GB/s modeled raw instance ceiling at 64 KiB operations.
Each volume is at least 48 GiB for the IOPS/capacity constraint. The service
model subsequently applies the 65\% EBS efficiency. Paying only average
throughput would not support the modeled active rebuild rate.

Reserve, per state and party, the persistent state $P$ plus
\[
 S=1.25\,[4N(2112+16t_N)]
\]
bytes of scratch: one largest tagged source array and 25\% working-space
allowance. Reserving scratch for all seven states permits suspended
rebuilds to retain work. This is a policy, not a proven maximum live-set
bound. The per-party volumes are striped/distributed logically; internal
record traffic is already charged rather than assuming direct shared disks.

For total workers $K_G+K_E=2(F+R)$, volumes $v$, capacity $D_v$ GiB,
throughput $T_v$ MiB/s, and IOPS $J_v$, the annual budget is
\[
 C_{\rm year}=63576.576(K_G+K_E)+
 \sum_v[0.96D_v+0.48(T_v-125)_++0.06(J_v-3000)_+].
\]
The baseline places both parties at private addresses in the same Availability Zone, so this
transfer placement has no separate EC2 regional data-transfer charge
\cite{aws-transfer}. Cross-AZ placement would need a different budget.

%%BUDGET%%

%%RAM%%

The annualization is steady-state at 10,000 queries/s, or 315.36 billion
queries/year. Initial database construction, cache warmup, client/application
transport, application computation, taxes, and replica-update consistency
are excluded. This is a semi-honest steady-state ledger; base-OT setup and
malicious-security checks are outside its scope. CPU is excluded by the
requested assumption. Each billed
fleet must still meet the advertised NIC/EBS rates in the actual placement;
both NIC cards and multiple flows are required for the modeled aggregate
network rate \cite{aws-network}.

\section{Query latency and batch size}
A central batcher receives 10,000 requests/s. For deterministic arrivals,
the mean batch-fill delay is $(B-1)/(2q)$: 51.15 ms at $B=1024$, 102.35 ms
at $B=2048$. It is not multiplied by the number of states. Splitting arrivals
into seven independent batchers would give a different fill delay.

At each dependent table stage, service is the maximum of NIC transfer,
EBS loading (including one 0.5 ms fixed wave allowance when needed), and
DRAM loading. Those per-stage maxima are summed. Taking one maximum over
the entire hierarchy would hide a level-local storage bottleneck. The
state/stash reads within a table are assumed parallel.

Read packed-value selection can wait until all metadata reads finish and
run as one batched OT layer. A parallel segmented propagation tree uses
at most $2\log_2B+1$ packed layers; inverse bitonic routing uses
$\log_2B(\log_2B+1)/2$. With $k=\log_2B$, the model's mean
RTT-equivalent dependency allowance is
\[
 U=2+A+(2k+1)+k(k+1)/2.
\]
Each unit costs 0.2 ms in the primary scenario. This is an explicit
scheduling allowance, not a proof of the exact compiled protocol's round
count. Preprocessed OT does not remove dependent online packed products.
The reported query latencies include fill, active service, and the finite
trace's admission queue.

%%LATENCYSENS%%

The effect of changing batch size is visible directly in the communication
ledger. Smaller batches save on the one-time bitonic work but add hierarchy
levels, rebuild frequency, and a larger fraction of fixed round/load delay.

%%BATCH%%

Reducing $B$ from 1024 to 512 saves only about 0.083 MB/query in the frontend;
total useful communication increases. In the resized five-state tests,
$B=512$ needs 59/78 workers per party and costs \$9.35M/\$13.11M annually,
with mean latency about 60--61 ms. It is not a cost-saving choice for the
stated tolerance. The headline seven-state $B=2048$ deployments save about
13--14\% versus $B=1024$ and retain sub-430 ms observed maxima in these
traces. Smaller batches remain relevant if a substantially tighter latency
target is imposed.

\clearpage
\section{Cross-cloud deployment}\label{sec:cross-cloud}
%%CROSSHEADLINE%%

AWS us-east-1 hosts the garbler and Google us-east4 hosts the evaluator,
connected privately in Northern Virginia. The workload remains a static
database with 32-byte keys and 32-byte values. Both scenarios retain the
construction, cache policy, and seven states; their RTTs are assumed.

%%CROSSSUMMARY%%

\subsection{Larger batches absorb round trips}
At $B=2048$, about 100 RTT-equivalent dependencies occupy the foreground
pool for half a second at 5 ms RTT, while a batch arrives every 0.205 seconds.
Relaxing response latency alone cannot make this pool sustain 10,000 reads/s.
More workers cannot remove the dependency delay. Keeping $B=2048$ would
require concurrent batch pipelines and a revised scheduler.

Larger batches give the pools time to finish each job. At $B=16{,}384$,
the arrival interval is 1.638 seconds and mean fill delay is 0.819 seconds,
supporting the 5 ms RTT scenario. At 10 ms RTT, $B=32{,}768$ gives a
3.277-second interval and 1.638-second mean fill delay. Removing lower
hierarchy levels also reduces communication per read despite the larger
frontend sort.

Cross-cloud pools target at most 90\% mean occupancy, including round and
loading waits, versus 70\% for same-AZ pools. This uses the extra latency
budget to reduce fleet and burst-bandwidth requirements. Two global
counter cycles include two largest-table rebuild starts per scenario;
reported latencies include the admission queue.

\subsection{Charge traffic in its actual direction}
About 98.6\% of cross-party bytes travel from garbler to evaluator.
AWS flat-rate Direct Connect covers that dominant direction; only the
smaller Google-to-AWS stream incurs a per-GiB transfer charge. Charge each
send once, excluding matching receives and same-party transfers.

%%CROSSTRAFFIC%%

\clearpage
\subsection{Private-link prices and capacity}
The price snapshot checked on 17 October 2026 combines an AWS Tier 1
100 Gbps Direct Connect port-pair at \$102.74/hour with two Google
100 Gbps Cross-Cloud Interconnect circuits at \$30/hour each and two
100 Gbps VLAN attachments at \$1/hour each
\cite{aws-flat-rate,gcp-interconnect}. Thus one redundant bundle costs
\[
 102.74+2(30+1)=\$164.74/{\rm hour},
\]
or \$1.443M/year. Its usable capacity is 100 Gbps, including a redundant
path; the two physical connections do not provide 200 Gbps of usable
capacity. AWS outbound transfer is included for the covered path. Google
charges \$0.02/GiB for outbound transfer on the selected US path and
nothing for inbound transfer. The selected topology connects primary
and redundant ports on both providers \cite{gcp-cci-topology}.

Average traffic alone understates the capacity needed by the scheduled
rebuild bursts. We reserve enough cross-cloud bandwidth for the foreground
and background pools to use their full modeled active NIC rates
simultaneously. For $F+R$ workers per party and $H$ usable 100 Gbps bundles,
the reservation satisfies
\[
 0.90\,H\,(12.5\ {\rm GB/s})\ \geq
 (F+R)(11.375\ {\rm GB/s}).
\]
The 10\% link allowance is separate from worker-pool occupancy. A second
check keeps mean traffic below 70\% of raw link capacity; the burst
condition determines the link counts in every scenario below.

%%CROSSCAPACITY%%

This reservation covers even the portion of worker NIC capacity used for
internal transfers, so cross-cloud bandwidth cannot further constrain the
modeled schedule. Those internal bytes do not enter the billed reverse
egress volume. With $q=10{,}000$ reads/s and $d_{\leftarrow}$ reverse
bytes/read, the annual networking cost is
\[
 C_{\rm net}=8760H(164.74)
 +\frac{(365\cdot24\cdot3600)q\,d_{\leftarrow}}{2^{30}}(0.02).
\]

Network charges dominate the total, so smoothing rebuild traffic across
time and sharing link capacity more tightly are useful optimization
targets. The reported budgets retain the full simultaneous-burst
reservation. A budget based only on average traffic would need a new
schedule before inheriting these latency estimates.

\paragraph{Scope of the estimate.}
Compute, storage, and I/O for both parties use the existing AWS prices as
a common infrastructure proxy. The CPU and storage assumptions are the
same as in the baseline. Provider selection must confirm the assumed RTT,
multi-terabit capacity, and eligibility of the chosen dedicated AWS
port-pair for flat-rate billing. Google workers and their VLAN attachments
must be in the same region to avoid extra regional transfer charges.
Any separately charged installation, cross-connect, colocation, or
optional network-service fees are outside this budget. The estimate is
a deployment projection from component prices, not a negotiated quote.

\clearpage
\section{Validation and remaining work}
All current resource paths consume the same materialized single-PRF case.
The annual model and report renderer reject a schedule stamped with a
different address protocol, preventing old traces from silently entering
new tables. Thirty-five tests, including six ORCompact/refinement checks,
cover direction conservation, private dummy
input prefetch in both directions, cache lifecycle, stagewise latency,
counter carries, headroom, price units, storage/state charging, label
translations, complete entry movement, and stale-schedule rejection.

The pruned AES computations match 32 independent native-AES key/plaintext
vectors for one through four projected bytes and ordinary prefixes, plus
a full AES known-answer test. The 32-AND S-box cost is held fixed so the
projection saving is isolated; newer S-box circuits are not silently
substituted. Functional pruning and ledger tests are not a compiled GC
benchmark or a security proof for the full composition.

The outstanding tasks are concrete: compile and validate the label/cache
interfaces and shared-bucket finalizer; implement the global planner's
external-memory schedule; measure directional NIC throughput, effective
DRAM loads, and EBS IOPS/waves; test real background preemption and bursty
arrivals; and extend the overflow/fingerprint failure accounting from
$2^{32}$ accesses to the annual workload. A year here is about 73.4 times
that horizon. The current numerical capacities must not be presented as
already proving an annual failure target.

Packing fingerprints only while moving them is a possible further
optimization, but its conversion/composition kernels and new I/O schedule
are not implemented. It is not included in any table in this report.
The priced construction uses Boolean fingerprints and packed values
consistently throughout.

\appendix
\section{Reproduction}
From the repository root, with Python, NumPy, and SciPy available:
\begin{verbatim}
python -m mpc_cost.hierarchical_throughput
python -m mpc_cost.hierarchical_communication_audit
python -m mpc_cost.hierarchical_prf_optimization
python -m mpc_cost.hierarchical_schedule
python -m mpc_cost.hierarchical_schedule --suite variants
python -m mpc_cost.hierarchical_annual
python -m mpc_cost.hierarchical_cross_cloud
python -m unittest discover -s mpc_cost -t . -p "test_hierarchical*.py"
python -m mpc_cost.render_hierarchical_throughput_report
\end{verbatim}
The scheduling commands rerun the full deterministic traces rather than
scaling an old price or latency by a byte-saving percentage. Current
case construction is in \texttt{hierarchical\_address.py}; decomposition,
trace, and price results are stored separately under \texttt{mpc\_cost/results}.
The original split ledger remains a reproducible input to the explicit
current-case transformation. It is not the protocol used for the final
resource tables. The standalone LaTeX includes all numerical tables and
needs no project-specific inputs to compile.

The cross-cloud command reruns the four larger-batch schedules and saves
their traffic, resource counts, link reservations, and costs in
\texttt{hierarchical\_cross\_cloud.json}. Published link prices and their
source URLs are recorded in
\texttt{mpc\_cost/data/cross\_cloud\_prices\_2026-10-17.json}.

\clearpage
\begin{thebibliography}{99}
\bibitem{construction} Microrama. User-supplied ongoing paper and the
project's Binary Hierarchical ORAM for Interactive Two-Party Computation
construction ledger, 17 October 2026.
\bibitem{orcompact} IACR ePrint 2022/1333. User-supplied ORCompact source,
Sections 3, 5.1, and Appendix B. \url{https://eprint.iacr.org/2022/1333}
\bibitem{nist-fips197} NIST. FIPS 197, Advanced Encryption Standard (AES),
updated May 2023. Sections 3.4 and 5.1.
\url{https://nvlpubs.nist.gov/nistpubs/FIPS/NIST.FIPS.197-upd1.pdf}
\bibitem{nist-aes} NIST. Circuit Complexity: List of Circuits. AES circuit
and S-box gate counts. Accessed 17 October 2026.
\url{https://csrc.nist.gov/projects/circuit-complexity/list-of-circuits}
\bibitem{nist-cmac} NIST. SP 800-38B: The CMAC Mode for Authentication.
\url{https://csrc.nist.gov/pubs/sp/800/38/b/upd1/final}
\bibitem{aws-compute} AWS. Specifications for EC2 compute optimized instances.
\url{https://docs.aws.amazon.com/ec2/latest/instancetypes/co.html}
\bibitem{aws-network} AWS. EC2 instance network bandwidth.
\url{https://docs.aws.amazon.com/AWSEC2/latest/UserGuide/ec2-instance-network-bandwidth.html}
\bibitem{aws-gp3} AWS. EBS General Purpose SSD volumes.
\url{https://docs.aws.amazon.com/ebs/latest/userguide/general-purpose.html}
\bibitem{aws-prices} AWS. AmazonEC2 us-east-1 public price list, version
20261015191316, published 15 October 2026. Selected rows and SKU identifiers
saved in \texttt{aws\_hierarchical\_prices\_2026-10-17.json}.
\url{https://pricing.us-east-1.amazonaws.com/offers/v1.0/aws/AmazonEC2/current/us-east-1/index.csv}
\bibitem{aws-transfer} AWS. Billing and metering for VPC owners and participants.
\url{https://docs.aws.amazon.com/vpc/latest/userguide/vpc-share-billing.html}
\bibitem{aws-flat-rate} AWS. Direct Connect flat-rate pricing.
Tier 1 100 Gbps port-pair, accessed 17 October 2026.
\url{https://aws.amazon.com/directconnect/pricing/flat-rate/}
\bibitem{gcp-interconnect} Google Cloud. Cloud Interconnect pricing,
Cross-Cloud Interconnect section. Accessed 17 October 2026.
\url{https://cloud.google.com/network-connectivity/interconnect/pricing}
\bibitem{gcp-cci-topology} Google Cloud. Cross-Cloud Interconnect overview.
\url{https://docs.cloud.google.com/network-connectivity/docs/interconnect/concepts/cci-overview}
\end{thebibliography}
\end{document}
